\section{Independent TT one-graviton mode-overlap normalization}
\label{app:tt-overlap}

The wave-zone coefficient $25/16$ can be obtained without using either the retarded source--receiver self-energy or a resonant absorption cross section. We summarize that independent route here.

Following the standard box-normalized TT quantization used in Ref.~\cite{BoughnRothman2006}, write
\begin{equation}
h_{ij}(\mathbf x,t)
=\frac{1}{c\sqrt V}
\sum_{\mathbf k,\lambda}
\sqrt{\frac{16\pi G\hbar}{\omega_k}}
\left[
a_{\mathbf k\lambda}
\epsilon^{(\lambda)}_{ij}
 e^{i(\mathbf k\cdot\mathbf x-\omega_k t)}
+\mathrm{h.c.}
\right],
\end{equation}
with $\epsilon^{(\lambda)}_{ij}\epsilon^{(\lambda')}_{ij}=\delta_{\lambda\lambda'}$. The local-inertial quadrupole interaction gives the one-graviton angular rate
\begin{equation}
\frac{d\kappa_g}{d\Omega}
=\frac{G\omega^5}{4\pi\hbar c^5}
\sum_\lambda
|Q_{ij}\epsilon^{(\lambda)}_{ij}|^2.
\label{eq:tt-diff-rate}
\end{equation}

For the plus transition
\begin{equation}
Q_{ij}=q\,\operatorname{diag}(1,-1,0),
\end{equation}
one finds
\begin{equation}
\sum_\lambda|Q:\epsilon^\lambda|^2
=q^2\mathcal F(\theta,\phi),
\end{equation}
where
\begin{equation}
\mathcal F
=2\cos^2\theta
+\frac12\sin^4\theta\cos^22\phi,
\end{equation}
and
\begin{equation}
\int d\Omega\,\mathcal F=\frac{16\pi}{5}.
\end{equation}
Equation~\eqref{eq:tt-diff-rate} therefore gives
\begin{equation}
\kappa_g
=\frac{4G\omega^5q^2}{5\hbar c^5},
\end{equation}
which is the linewidth normalization used in the main text.

Define the normalized one-graviton angular mode
\begin{equation}
\boxed{u_\lambda(\hat{\mathbf n})
=\sqrt{\frac{5}{16\pi}}
\frac{Q:\epsilon^\lambda(\hat{\mathbf n})}{q},}
\end{equation}
so that $\sum_\lambda\int d\Omega|u_\lambda|^2=1$. For an identical aligned reciprocal receiver displaced by $R\hat z$, a plane-wave component acquires the phase $e^{iz\mu}$ with $z=kR$ and $\mu=\cos\theta$. The normalized fixed-frequency mode overlap is therefore
\begin{equation}
S(z)
=\sum_\lambda\int d\Omega
|u_\lambda|^2e^{iz\mu}
=\frac{5}{32}\int_{-1}^{1}d\mu\,
(1+6\mu^2+\mu^4)e^{iz\mu}.
\label{eq:tt-overlap-integral}
\end{equation}
Direct evaluation gives
\begin{equation}
S(z)
=\frac{5}{4z^5}
\left[
2z^4\sin z+4z^3\cos z
-6z^2\sin z-6z\cos z+6\sin z
\right].
\end{equation}
Writing the result as positive- and negative-time propagation pieces,
\begin{equation}
S(z)=S_+(z)+S_-(z),
\end{equation}
with
\begin{equation}
\boxed{
S_+(z)
=-\frac{5i}{4}
\frac{P(z)e^{iz}}{z^5},}
\end{equation}
\begin{equation}
S_-(z)
=+\frac{5i}{4}
\frac{P(-z)e^{-iz}}{z^5},
\end{equation}
where
\begin{equation}
P(z)=3-3iz-3z^2+2iz^3+z^4.
\end{equation}
For a wavepacket carrying the usual $e^{-i\omega t}$ dependence, $S_+$ has phase $e^{-i\omega(t-R/c)}$ and is the outgoing retarded source-to-receiver component; $S_-$ is the time-reversed component of the full fixed-frequency overlap. Thus
\begin{equation}
\boxed{
t_{BA}^{\rm TT}(z)
=-\frac{5i}{4}
\frac{P(z)e^{iz}}{z^5}.}
\label{eq:tt-exact-transfer}
\end{equation}
The polynomial in Eq.~\eqref{eq:tt-exact-transfer} emerges from the TT mode overlap and is the same polynomial obtained independently from the retarded electric-Weyl calculation.

In the wave zone,
\begin{equation}
t_{BA}^{\rm TT}
=-\frac{5i}{4}\frac{e^{iz}}{z}
\left[1+O(z^{-1})\right],
\end{equation}
so
\begin{equation}
\boxed{
\eta_{\rm store}
=|t_{BA}^{\rm TT}|^2
\longrightarrow
\frac{25}{16(kR)^2}.}
\end{equation}
This is the third independent normalization route used to validate the propagation coefficient: canonical TT one-graviton mode overlap, in addition to the retarded-field and critical-coupling/Friis derivations.

The separated $S_+$ component should not be interpreted as a bounded single-pass probability in the reactive near zone; the storage interpretation is the weak one-way wave-zone regime used in the link budget.
